\documentclass[11pt]{article}
\usepackage[margin=1in]{geometry}
\usepackage{setspace,microtype}
\usepackage[T1]{fontenc}
\usepackage{newtxtext,newtxmath}
\usepackage{siunitx}
\setstretch{1.15}
\setlength{\parindent}{0pt}
\setlength{\parskip}{0.55em}
\pagestyle{empty}
\newcommand{\PubValidationCases}{40}
\newcommand{\PubValidationMaxError}{\num{5.59e-08}}
\newcommand{\PubValidationMinSlack}{\num{0.00e+00}}
\newcommand{\PubValidationCertViolation}{\num{0.00e+00}}
\newcommand{\PubValidationCertGap}{\num{1.00e-08}}
\newcommand{\PubKernelCases}{48}
\newcommand{\PubKernelLargeSpeedup}{21.05}
\newcommand{\PubKernelMaxError}{\num{9.78e-09}}
\newcommand{\PubKernelCompressedSlope}{0.94}
\newcommand{\PubKernelDenseSlope}{2.00}
\newcommand{\PubTraceCases}{24}
\newcommand{\PubTraceIntervals}{168}
\newcommand{\PubTraceGeoSpeedup}{2.53}
\newcommand{\PubTraceWins}{24}
\newcommand{\PubTraceDominancePct}{100.0\%}
\newcommand{\PubTraceCertGap}{\num{9.56e-09}}
\newcommand{\PubPrimaryCases}{60}
\newcommand{\PubPrimaryGeoSpeedup}{2.33}
\newcommand{\PubPrimaryCILow}{2.16}
\newcommand{\PubPrimaryCIHigh}{2.51}
\newcommand{\PubPrimaryWins}{60}
\newcommand{\PubPrimaryMaxDifference}{\num{0.00e+00}}
\newcommand{\PubPrimaryCertGap}{\num{9.43e-09}}
\newcommand{\PubPrimaryMedianThetaPct}{28.8\%}
\newcommand{\PubPrimarySignP}{\num{8.67e-19}}
\newcommand{\PubPrimaryNLowSpeedup}{1.61}
\newcommand{\PubPrimaryNHighSpeedup}{3.35}
\newcommand{\PubRepeatBlocks}{300}
\newcommand{\PubRepeatWins}{298}
\newcommand{\PubRepeatMinSpeedup}{0.79}
\newcommand{\PubCompressedMedianCV}{0.60\%}
\newcommand{\PubCliqueMedianCV}{0.39\%}
\newcommand{\PubCliqueMedianNnz}{\num{100394}}
\newcommand{\PubCliqueMedianCSRMiB}{1.19}
\newcommand{\PubFamilyCorrelated}{2.33}
\newcommand{\PubFamilyDense}{6.30}
\newcommand{\PubFamilyBreakpoints}{1.34}
\newcommand{\PubFamilyNearTie}{1.76}
\newcommand{\PubRobustnessCases}{36}
\newcommand{\PubStressCases}{8}
\newcommand{\PubStressGeoSpeedup}{5.80}
\newcommand{\PubStressNHighSpeedup}{6.71}
\newcommand{\PubApplicationCases}{9}
\newcommand{\PubApplicationWins}{9}
\newcommand{\PubApplicationGeoSpeedup}{1.46}
\newcommand{\PubApplicationCILow}{1.24}
\newcommand{\PubApplicationCIHigh}{1.75}
\newcommand{\PubApplicationNHighSpeedup}{1.70}
\newcommand{\PubExternalCases}{9}
\newcommand{\PubExternalWins}{9}
\newcommand{\PubExternalGeoSpeedup}{2.31}
\newcommand{\PubExternalCILow}{2.22}
\newcommand{\PubExternalCIHigh}{2.40}
\newcommand{\PubExternalNHighSpeedup}{3.07}
\newcommand{\PubExternalCertGap}{\num{9.85e-09}}


\begin{document}

\begin{center}
{\large\bfseries Executive Summary}\\[0.35em]
{\normalsize Simultaneous Group-Envelope Bounds for $\Gamma$-Robust Multiple-Choice Knapsack Problems}\\[0.35em]
{\small August 2026}
\end{center}

Many resource-allocation decisions choose exactly one option per group.  Protecting their shared constraint with a Bertsimas--Sim uncertainty budget produces a robust multiple-choice knapsack problem.  The classical threshold reduction is polynomial, but many distinct deviations can make complete LP certification expensive.  B\"using, Gersing, and Koster (2023) report that LP solves consume 33.1\% of runtime on average among solved instances and 93.17\% at one million items, where the root LP alone averages 1,062 seconds.  Their objective-uncertainty setting differs from ours, but it identifies the bounded-threshold LP family as a material bottleneck.

After dualizing a fixed-threshold MCKP, all threshold-dependent baselines cancel.  Between consecutive group deviations, the remaining envelope is the maximum of one constant and one common-slope line.  Prefix and suffix maxima with range accumulation evaluate one multiplier over all $B$ thresholds in $O(B+K\log(B+1))$ time and $O(B+K)$ working storage, rather than $\Theta(BK)$ direct score evaluations, for $K$ options.

Weak duality makes the interval minimax value a valid upper bound on every feasible fixed-threshold LP; strong LP duality gives equality on feasible singletons.  An epigraph-dual mapping proves exact-minimax dominance over the bounded-threshold group-clique LP.  Convex bracketing returns a valid bound with an explicit minimization-error certificate, and best-first splitting certifies the maximum LP value over the complete threshold family.

The comparison is component-matched: both certificates use the same adaptive search, fixed-threshold solver, candidate rule, tolerance, and timing policy; only the interval upper bound differs.  A common-trace experiment also evaluates both oracles on identical intervals.  This isolates the new component rather than comparing with DnC+, a full Java/Gurobi solver that combines filtering, estimators, cuts, incumbent operations, and early termination for objective uncertainty.

Across \PubPrimaryCases{} primary instances, the compressed certificate reaches tolerance throughout, wins all \PubPrimaryWins{} instance-level median comparisons, and is \PubPrimaryGeoSpeedup{} times faster in geometric mean than the sparse clique LP (95\% interval \PubPrimaryCILow--\PubPrimaryCIHigh).  Speedup rises from \PubPrimaryNLowSpeedup{} at 360 groups to \PubPrimaryNHighSpeedup{} at 1,440 groups.  All 36 robustness cases reach tolerance.  On \PubExternalCases{} published robust-knapsack coefficient sets, it wins all \PubExternalWins{} comparisons with a \PubExternalGeoSpeedup-fold advantage; this is an out-of-generator coefficient test, not replication of the source model.  Algebraic, epigraph-LP, and exhaustive-scan checks find maximum error \PubValidationMaxError{} and no certificate violation.

The deliverable is a faster complete threshold-disjunctive LP certificate for use at the root or in an outer robust search.  A separate exact-integer audit confirms objective agreement and valid gap accounting, while showing that integer subproblem work can erase the oracle advantage and that compact SCIP remains preferable on many-breakpoints instances.  The claim is therefore simultaneous all-threshold evaluation for exactly-one groups, not a new uncertainty model or universal integer-solver superiority.

\end{document}
